% tex8web: text, cross-references, theorems, the table of contents, lists and
% figures. (tcolorbox's boxes are in tex8web-boxes.tex)
\makeatletter

% --- title ------------------------------------------------------------------

% a \thanks's mark where it is, as in the PDF: tex4ht (latex.4ht) puts it
% there only where \footnote is \thanks, which LaTeX's \maketitle makes it
% but titling's doesn't; and the marks not links, as in the PDF
\let\web@thanks\thanks
\def\thanks#1{{\let\footnote\thanks\web@thanks{#1}}}
\Configure{thank}
    {\csname a:thank-mark\endcsname}{\csname b:thank-mark\endcsname}
    {\HCode{<p>}\csname a:thank-mark\endcsname}{\csname b:thank-mark\endcsname}
    {\HCode{</p>}}

% --- links ------------------------------------------------------------------

\@ifpackageloaded{hyperref}{%
    % tex4ht's \nolinkurl chokes on "_" inside \author. (A # comes out
    % doubled)
    \renewcommand\nolinkurl[1]{\texttt{\detokenize{#1}}}%
    % tex4ht writes the table of contents with \texorpdfstring's plain-text
    % argument, meant for PDF bookmarks; use the TeX one, as the PDF's table
    % of contents does
    \Configure{writetoc}{\let\texorpdfstring\@firstoftwo}%
}{}

% tex4ht's cleveref links only the number of \cref; hand cleveref the link
% start and end instead, so that with nameinlink all of "Lemma 1" is the
% link, as in the PDF
\@ifpackageloaded{cleveref}{%
    \def\web@setcref#1#2{\bgroup
        \def\rEfLiNK##1##2{\Link{##1}{}}%
        \def\XRrEfLiNK[##1]##2##3{\Link[##1]{##2}{}}%
        \cref@getlabel{#2}{\@templabel}%
        #1{\@templabel}%
          {\expandafter\expandafter\expandafter\@firstoffive\csname r@#2\endcsname{}{}{}{}{}}%
          {\EndLink}%
        \egroup
    }%
    \HLet\@@@setcref\web@setcref
}{}

% --- theorems ---------------------------------------------------------------

% spaced as in the PDF rather than tex4ht's 2em: amsthm's theorem style's
% space above and below, or LaTeX's \topsep. The style sets the space only
% once the theorem has started, so leave it in the heading, for
% text.theorem_spacing to move onto the theorem (\HCode would expand it at
% shipout, too late)
\ifdefined\thm@preskip
    \def\web@thmskip{\web@em\thm@preskip}%
    \def\web@thmpostskip{\web@em\thm@postskip}%
\else
    \def\web@thmskip{\web@em\topsep}%
    \let\web@thmpostskip\web@thmskip
\fi
\Configure{newtheorem}
    {\ifvmode\IgnorePar\fi\EndP
     \HCode{<div class="newtheorem">}\par\noindent\HCode{<span class="head">}}
    {\edef\web@mark{\noexpand\HCode{<span class="tex8web-thmskip"
         data-before="\web@thmskip" data-after="\web@thmpostskip"></span>}}%
     \web@mark\HCode{</span>}}
    {\ifvmode\IgnorePar\fi\EndP\HCode{</div>}\par\indent\ShowPar}

% --- table of contents ------------------------------------------------------

% section numbers in a box as wide as in the PDF (tocloft's widths, or the
% standard classes'), so titles line up whether a number has one digit or
% two, and the space above each entry as in the PDF: tocloft's, or
% article's (the third argument would add page numbers; see
% css/tex8web.css)
\ConfigureToc{section}
    {\HCode{<span class="sectionToc"><span class="tocnumber">}}
    {\HCode{</span>}}
    {}
    {\HCode{</span><br />}}
\ConfigureToc{subsection}
    {\HCode{<span class="subsectionToc"><span class="tocnumber">}}
    {\HCode{</span>}}
    {}
    {\HCode{</span><br />}}
\ifdefined\cftsecnumwidth
    \web@Css{:root { --tex8web-toc-number: \web@em\cftsecnumwidth;
        --tex8web-toc-indent: \web@em\cftsubsecindent;
        --tex8web-toc-subnumber: \web@em\cftsubsecnumwidth;
        --tex8web-toc-section-skip: \web@em\cftbeforesecskip;
        --tex8web-toc-subsection-skip: \web@em\cftbeforesubsecskip; }}
\else
    \web@Css{:root { --tex8web-toc-number: \web@em{1.5em};
        --tex8web-toc-indent: \web@em{1.5em};
        --tex8web-toc-subnumber: \web@em{2.3em};
        --tex8web-toc-section-skip: \web@em{1em};
        --tex8web-toc-subsection-skip: 0em; }}
    % and section entries in bold. (tocloft's own font commands aren't
    % followed)
    \Css{.sectionToc, .likesectionToc { font-weight: bold; }}
\fi

% a table of contents a document makes with \@starttoc{toc} itself (in a box,
% say) is tex4ht's, as \tableofcontents makes without its heading. (tex4ht
% writes the entries in its own form, which only it reads)
\let\web@starttoc\@starttoc
\def\@starttoc#1{%
    \def\web@a{#1}\def\web@b{toc}%
    \ifx\web@a\web@b
        \expandafter\@firstoftwo
    \else
        \expandafter\@secondoftwo
    \fi
    {\csname :tableofcontents\endcsname}{\web@starttoc{#1}}}

% --- lists ------------------------------------------------------------------

% LaTeX's spacing of a list: \topsep around it, \itemsep and \parsep between
% items (see css/tex8web.css)
\begingroup\@listi
    \web@Css{:root { --tex8web-list-sep: \web@em\topsep;
        --tex8web-item-sep: \web@em{\dimexpr\itemsep+\parsep\relax}; }}
\endgroup

% enumitem's enumerate, and its lists made with \newlist: html4.4ht's grid of
% labels and items, with each list's own lengths, set by the first \item,
% where tex4ht opens the list (latex.4ht). (tex4ht's names need : a letter)
\catcode58 11
\@ifpackageloaded{enumitem}{%
    \def\web@enumitem#1#2{%
        \ConfigureList{#1}
            {\EndP\edef\web@list{\noexpand\HCode{<dl \noexpand\a:LRdir
                 class="enumerate-enumitem#2"
                 style="--tex8web-enumitem-topsep: \web@em\topsep;
                 --tex8web-enumitem-itemsep: \web@em{\dimexpr\itemsep+\parsep\relax};
                 --tex8web-enumitem-indent: \web@em\leftmargin;
                 --tex8web-enumitem-labelsep: \web@em\labelsep">}}\web@list
             \PushMacro\end:itm \global\let\end:itm=\empty}
            {\PopMacro\end:itm \global\let\end:itm \end:itm
             \EndP\HCode{</dd></dl>}\ShowPar}
            {\end:itm \global\def\end:itm{\EndP\Tg</dd>}\HCode{<dt
                 class="enumerate-enumitem#2">}\bgroup \bf}
            {\egroup\EndP\HCode{</dt><dd\Hnewline class="enumerate-enumitem">}}}%
    \web@enumitem{enumerate}{}%
    % (html4.4ht's own names the list last made, not this one)
    \def\web@newlist#1{\web@enumitem{#1}{ enumitem-#1}}%
    \Configure{enumitemnewlist}{\expandafter\web@newlist\expandafter{\a:enumitemenv}}%
}{}
\catcode58 12

% --- figures ----------------------------------------------------------------

% \includegraphics[width=..] becomes a percentage of the column, as in the
% PDF; other sizes as tex4ht writes them
\ExplSyntaxOn
\cs_if_exist:NT \Gin@req@width
  {
    \cs_set_eq:Nc \__web_gin_dim: { a:Gin-dim }
    \Configure{Gin-dim}{%
        \expandafter\ifx\csname Gin:ewidth\endcsname\Gin@exclamation
            \__web_gin_dim:
        \else
            style="width:\fp_eval:n{round(\Gin@req@width/\linewidth*100,2)}\char_generate:nn{`\%}{12}"
        \fi
    }
  }
\ExplSyntaxOff

\makeatother
